% Appendices for the MEDORA system-analysis and design report.
% All API contracts and test cases below are specifications; none is execution evidence.

\section{Detailed Use Case Specifications}
The following specifications expand the highest-risk use cases from Section 3. They are the behavior baseline for implementation and verification. \weak{Facility policies, clinical rules, approval thresholds and retention requirements must be approved before the matching behavior is final.}

\subsection{UC-021 Book Appointment}
\begin{description}[leftmargin=2.6cm,style=nextline]
\item[Goal] Reserve a valid clinic service slot for an identified patient.
\item[Primary actor] Receptionist. Patient self-service is Phase 2.
\item[Trigger] A patient requests an appointment or a receptionist schedules follow-up.
\item[Preconditions] Actor is authenticated and facility-scoped; patient identity and requested service are resolved; scheduling data is available.
\item[Main success flow]
\begin{enumerate}[leftmargin=1.4em]
  \item Receptionist selects patient, service, provider, and requested time.
  \item System validates request fields, actor permissions, facility, and patient scope.
  \item System checks provider schedule, duration, room, and any configured constrained resource.
  \item System repeats the availability check at transaction commit and atomically creates the Appointment and required ResourceBooking records.
  \item System records actor/time and returns appointment ID and status.
\end{enumerate}
\item[Alternative flows] If a slot is unavailable, the system offers permitted alternatives; the receptionist selects one and the system repeats validation. A draft request may be retained without reserving resources if the clinic enables drafts.
\item[Exception flows] Invalid identity, unauthorized scope, unavailable resource, or a concurrent conflict rejects the request without partial appointment/resource records. Repeating a successful request with the same idempotency key returns the original result.
\item[Postconditions and data] Success creates one Appointment and optional resource allocation; it does not create an Encounter. Failure leaves the prior schedule unchanged.
\item[Authorization] Receptionist in facility scope; any patient portal actor is separately authenticated and limited to their own identity.
\item[Trace] BR-002; FR-005, FR-015, FR-017; UC-020--UC-026; TC-003; Figures~\ref{fig:sequence-book-appointment} and~\ref{fig:state-appointment}.
\end{description}

\subsection{UC-027 Check In Patient}
\begin{description}[leftmargin=2.6cm,style=nextline]
\item[Goal] Start one valid outpatient Encounter and place it in the configured queue.
\item[Primary actor] Receptionist.
\item[Trigger] An identified booked patient arrives or a walk-in is accepted by facility policy.
\item[Preconditions] Receptionist is authenticated and scoped; patient identity is resolved; walk-in policy is known where applicable.
\item[Main success flow]
\begin{enumerate}[leftmargin=1.4em]
  \item Receptionist selects an Appointment or accepted walk-in intake.
  \item System validates patient identity, facility scope, and current Appointment state.
  \item System checks whether the patient already has an active Encounter under the approved rule.
  \item System creates/activates the Encounter and assigns the configured queue state atomically.
  \item System returns Encounter and queue references.
\end{enumerate}
\item[Alternative flows] Accepted walk-ins use the same identity and duplicate-Encounter checks. A rejected walk-in remains unregistered as an active Encounter and receives the permitted next step.
\item[Exception flows] Unresolved identity, unauthorized facility, invalid Appointment state, or a duplicate active Encounter rejects the transition without creating a second Encounter.
\item[Postconditions and data] Success links Encounter to one Patient and optionally one Appointment. A booking without check-in remains an Appointment only.
\item[Authorization] Receptionist may create/check in; Nurse may read only the queue assigned by role/facility policy.
\item[Trace] BR-003; FR-006; UC-027--UC-030; TC-004; Figures~\ref{fig:sequence-check-in} and~\ref{fig:state-encounter}.
\end{description}

\subsection{UC-037/UC-038 Sign and Amend Clinical Note}
\begin{description}[leftmargin=2.6cm,style=nextline]
\item[Goal] Finalize a clinically attributable note while preserving any later correction history.
\item[Primary actor] Doctor assigned to the active Encounter.
\item[Trigger] Doctor completes documentation or identifies an error in signed content.
\item[Preconditions] Doctor is authenticated, authorized for the Encounter, and the encounter is in a state that permits documentation.
\item[Main success flow]
\begin{enumerate}[leftmargin=1.4em]
  \item Doctor saves a draft with the required clinical sections.
  \item System validates required data, assignment, and expected current version.
  \item Doctor requests signing.
  \item System records an immutable signed version with actor and timestamp and emits audit evidence.
  \item System returns the signed-version reference.
\end{enumerate}
\item[Alternative flows] An incomplete note remains a Draft. A correction creates a successor amendment linked to the signed predecessor; the Doctor supplies a reason and signs the amendment separately.
\item[Exception flows] Missing fields, stale version, invalid assignment, or unauthorized amendment rejects the operation. The signed source and all previous versions remain unchanged.
\item[Postconditions and data] NoteVersion history is append-only; no direct update of signed content is allowed.
\item[Authorization] Assigned Doctor may sign/amend; all reads/writes remain role-, facility-, and encounter-scoped.
\item[Trace] BR-005; FR-009, FR-020, FR-032; UC-037/UC-038; TC-005; Figures~\ref{fig:sequence-clinical-note} and~\ref{fig:state-clinical-note}.
\end{description}

\subsection{UC-048 Validate and Issue Prescription}
\begin{description}[leftmargin=2.6cm,style=nextline]
\item[Goal] Issue an authorized, complete prescription after configured safety checks.
\item[Primary actor] Doctor assigned to the active Encounter.
\item[Trigger] Doctor submits a prescription draft for issue.
\item[Preconditions] Patient and Encounter are verified; Doctor is assigned and authorized; at least one medicine item is present.
\item[Main success flow]
\begin{enumerate}[leftmargin=1.4em]
  \item Doctor submits medicine items and directions.
  \item System validates required fields and configured rule inputs that are available.
  \item System presents any supported warning and its reason.
  \item Doctor resolves the warning or performs a permitted override with reason.
  \item System issues the Prescription, records actor/time/audit evidence, and returns the issue reference.
\end{enumerate}
\item[Alternative flows] Doctor may keep the Prescription in Draft. If no blocking warning is raised, the Doctor may issue after required-field validation.
\item[Exception flows] Missing data, a blocking warning, unavailable required policy data, unauthorized override, or stale Encounter state prevents issue and creates no administration event.
\item[Postconditions and data] Issued Prescription and PrescriptionItems are distinct from MedicationAdministration; issuing alone never asserts medicine was dispensed/administered.
\item[Authorization] Assigned Doctor issues; medicine staff record only the separately authorized administration/dispensing event.
\item[Trace] BR-006/BR-008; FR-011--FR-013; UC-043--UC-052; TC-006; Figure~\ref{fig:sequence-prescription}.
\end{description}

\subsection{UC-039/UC-040 Schedule and Execute Service Order}
\begin{description}[leftmargin=2.6cm,style=nextline]
\item[Goal] Coordinate a requested service with valid resources and record its actual outcome.
\item[Primary actors] Doctor creates the order; assigned Service Staff records execution.
\item[Trigger] Doctor adds a service/procedure to an active Encounter.
\item[Preconditions] The Encounter is active; service catalog and resource requirements are configured; actors have relevant scope.
\item[Main success flow]
\begin{enumerate}[leftmargin=1.4em]
  \item Doctor records service, reason, priority, and required preparation.
  \item System validates Encounter, order fields, and actor permissions.
  \item Scheduler requests a time/resource allocation.
  \item System validates room, staff, and constrained-resource availability at commit.
  \item Assigned staff starts the order and records actor/time.
  \item Staff records required outcome; system validates and marks execution complete.
\end{enumerate}
\item[Alternative flows] If no resource slot is available, the order remains pending without a booking and may be rescheduled after a fresh availability check. A delayed order remains non-complete.
\item[Exception flows] Conflict, missing assignment, invalid state, or incomplete outcome rejects the transition and preserves the order's last valid state. Cancellation records a reason and releases resources according to policy.
\item[Postconditions and data] ProcedureOrder, ResourceBooking, and execution outcome are linked but distinct records.
\item[Authorization] Doctor creates order; assigned Service Staff updates execution; scheduling authority follows clinic policy.
\item[Trace] BR-002/BR-007; FR-014--FR-018; UC-039/UC-040; TC-007; Figures~\ref{fig:sequence-service-order} and~\ref{fig:activity-service-order}.
\end{description}

\subsection{UC-063/UC-065 Record Payment and Refund}
\begin{description}[leftmargin=2.6cm,style=nextline]
\item[Goal] Record settlement or refund while preserving a complete financial history.
\item[Primary actors] Cashier posts payment; authorized Finance Staff posts refund/adjustment.
\item[Trigger] Eligible Invoice is issued or a prior PaymentTransaction requires refund/adjustment.
\item[Preconditions] Actor is authenticated and facility/finance-scoped; invoice and original transaction state are known.
\item[Main success flow]
\begin{enumerate}[leftmargin=1.4em]
  \item Cashier selects invoice and enters amount, currency, and method.
  \item System validates balance, state, scope, and idempotency key.
  \item System appends PaymentTransaction and returns the resulting invoice balance.
  \item For a refund, Finance Staff references the original transaction, supplies a reason, and obtains configured approval.
  \item System appends a refund/adjustment transaction and audit event; the original payment remains intact.
\end{enumerate}
\item[Alternative flows] Partial payment leaves an outstanding balance. Same-key retries return the original result.
\item[Exception flows] Invalid amount/state, conflicting idempotency key, insufficient authority, or missing reason/approval has no financial effect.
\item[Postconditions and data] Invoice and PaymentTransaction history is preserved; clinical records are unchanged.
\item[Authorization] Cashier posts payments; a separately authorized finance role handles refunds and approvals.
\item[Trace] BR-009; FR-022/FR-023; UC-063/UC-065; TC-008/TC-009; Figure~\ref{fig:sequence-payment-refund}.
\end{description}

\section{Supplementary UML Diagrams}
Table~\ref{tab:uml-artifact-inventory} catalogs the editable sources, rendered figures and traceability inventory. \weak{The diagrams are models, not implementation evidence.} Section 4 is reserved for the architecture, component, class and ERD diagrams; workflow and lifecycle diagrams are in Section 5.

\begin{longtable}{@{}p{0.26\textwidth}p{0.16\textwidth}p{0.42\textwidth}p{0.09\textwidth}@{}}
\caption{UML source and rendered-figure inventory}\label{tab:uml-artifact-inventory}\\
\toprule Source group & Types / count & Purpose & Status \\
\midrule
\endfirsthead
\toprule Source group & Types / count & Purpose & Status \\
\midrule
\endhead
Section 2 business analysis & Activity 2; domain 1 & Outpatient journey, appointment/check-in exceptions, and clinical domain concepts. & Rendered \\
Section 3 requirements & Functional 1; overall use case 1; focused use case 9 & Capability groups, system boundary, actors, and detailed actor goals. & Rendered \\
Section 4 architecture & Reserved for Chapter 4 owner & This contribution restores the original Chapter 4 outline and claims no Chapter 4 UML. & \weak{Pending} \\
Section 4 persistence & Reserved for Chapter 4 owner & ERD and schema artifacts remain with the Chapter 4 contributor. & \weak{Pending} \\
Section 4 classes & Reserved for Chapter 4 owner & Structural class artifacts remain with the Chapter 4 contributor. & \weak{Pending} \\
Section 5 interaction & System sequence 6 & Actor-visible success, alternative, and rejection outcomes for high-risk use cases. & Rendered \\
Section 5 process/lifecycle & Activity 2; state machine 5 & Service execution, closure/finance separation, and entity transitions. & Rendered \\
Section 5 UI wireframes & Wireframe 3 & Front desk and queue, clinician encounter, and cashier invoice and payment screens. Captures of a running application come after implementation. & Rendered \\
\bottomrule
\end{longtable}

The PlantUML source is stored in the matching \code{uml_code/} folder and the rendered PNG in \code{figure/}. Rendering checks syntax and supports visual review; \weak{it does not establish implementation, clinical approval or user acceptance.}

\section{API Specifications and Data Schemas}
This appendix and Chapter 5 propose API request and response examples and a common error schema. The examples below show representative persisted data shapes for implementation planning; all IDs and values are synthetic.

\subsection{Representative Appointment Data}
\begin{lstlisting}[language=json,caption={Proposed Appointment resource shape},label={lst:appointment-resource}]
{
  "appointmentId": "appt-demo-1001",
  "patientId": "pat-demo-001",
  "serviceId": "svc-demo-012",
  "providerId": "emp-demo-004",
  "facilityId": "facility-demo-main",
  "startAt": "2026-10-20T09:30:00+07:00",
  "endAt": "2026-10-20T10:00:00+07:00",
  "status": "BOOKED",
  "createdAt": "2026-10-16T09:00:00+07:00"
}
\end{lstlisting}

\subsection{Representative Encounter Data}
\begin{lstlisting}[language=json,caption={Proposed Encounter resource shape},label={lst:encounter-resource}]
{
  "encounterId": "enc-demo-2001",
  "patientId": "pat-demo-001",
  "appointmentId": "appt-demo-1001",
  "status": "WAITING",
  "facilityId": "facility-demo-main",
  "assignedClinicianId": "emp-demo-004",
  "startedAt": "2026-10-20T09:25:00+07:00"
}
\end{lstlisting}

\subsection{Representative Audit Event Data}
\begin{lstlisting}[language=json,caption={Proposed audit event shape with no clinical payload duplication},label={lst:audit-event-resource}]
{
  "auditEventId": "audit-demo-9001",
  "actorId": "user-demo-001",
  "action": "CLINICAL_NOTE_SIGNED",
  "targetType": "ClinicalNoteVersion",
  "targetId": "note-version-demo-002",
  "facilityId": "facility-demo-main",
  "occurredAt": "2026-10-20T10:02:00+07:00",
  "reason": null,
  "correlationId": "req-demo-001"
}
\end{lstlisting}

These shapes are illustrative contract baselines. \weak{The approved OpenAPI schema, database types, retention rules, field-level access policy and migration definitions still have to be written and reviewed before implementation.}

\section{Detailed Test Cases and Execution Results}
Detailed test procedures and expected outcomes are listed in Table~\ref{tab:detailed-test-procedures}. The actual-result, tester, timestamp, build and evidence fields are left empty. \weak{All 13 cases remain \textbf{Not Run} because no application or test environment exists, and no pass or fail result is claimed.}

\begingroup
\scriptsize
\renewcommand{\arraystretch}{1.08}
\begin{longtable}{@{}p{0.06\textwidth}p{0.22\textwidth}p{0.27\textwidth}p{0.31\textwidth}p{0.08\textwidth}@{}}
\caption{Detailed test procedures and expected outcomes}\label{tab:detailed-test-procedures}\\
\toprule ID & Setup & Procedure & Expected result / evidence to capture & Status \\
\midrule
\endfirsthead
\toprule ID & Setup & Procedure & Expected result / evidence to capture & Status \\
\midrule
\endhead
TC-001 & Synthetic patient; authenticated user without clinical permission. & Call protected patient/clinical read and write endpoints directly, bypassing the UI. & Reject both requests; disclose no protected fields; create no mutation; retain sanitized response and audit-policy evidence. & \weak{Not Run} \\
TC-002 & One synthetic patient; registrar and ordinary receptionist accounts. & Register a matching identity, then request merge as receptionist. & Flag possible duplicate; deny unauthorized merge; preserve source records and history. & \weak{Not Run} \\
TC-003 & One provider/room/resource slot; two concurrent clients. & Submit conflicting bookings with separate idempotency keys at the same time. & At most one booking commits; losing request receives conflict/alternatives; no orphan resource reservation remains. & \weak{Not Run} \\
TC-004 & Patient with no active encounter and an accepted appointment or walk-in. & Check in once; repeat check-in while the first encounter remains active. & First call creates one encounter/queue assignment; second is rejected without duplicate state. & \weak{Not Run} \\
TC-005 & Assigned doctor; incomplete draft and then corrected note. & Attempt signing while required content is missing; correct, sign, then append an amendment with reason. & Incomplete signing is blocked; signed source remains immutable; amendment links predecessor and records actor/time/reason. & \weak{Not Run} \\
TC-006 & Synthetic, clinician-approved safety-rule fixture; doctor and medicine staff roles. & Issue with a warning; test dismissal and permitted override; inspect medication events. & Apply configured warning behavior; require authorized override reason; prescription issue alone creates no administration event. & \weak{Not Run} \\
TC-007 & Valid order and two requests competing for the same resource. & Attempt conflicting allocations; submit incomplete execution result and then a valid result. & Conflict does not commit; incomplete result cannot complete order; valid result records performer/time and releases resources by policy. & \weak{Not Run} \\
TC-008 & Issued invoice with known total; cashier in authorized facility. & Post partial payment; retry same key; attempt another payment with a new key. & Same key returns original event; permitted new payment updates balance exactly once; reconciliation matches transaction history. & \weak{Not Run} \\
TC-009 & Paid synthetic invoice; cashier and authorized finance role. & Request refund with unauthorized role, invalid amount, missing reason, then valid approval. & Invalid requests have no financial effect; valid refund appends event and preserves original payment. & \weak{Not Run} \\
TC-010 & Two facility scopes with synthetic report data. & Query and export through auditor/manager roles across authorized and unauthorized scopes. & Enforce scope; audit actor, time, filter, and target; do not duplicate protected clinical content in logs. & \weak{Not Run} \\
TC-011 & Staging-like environment, agreed synthetic dataset, representative concurrent load. & Measure ordinary screens and search; record sample size, p95, build, and configuration. & Compare results with owner-approved targets; retain sanitized measurements; do not infer production capacity. & \weak{Not Run} \\
TC-012 & Encrypted synthetic backup and clean recovery environment. & Restore database and private files; verify checksums and measure data loss and recovery time. & Compare measured RPO/RTO with approved objectives; retain restore log and timestamps. & \weak{Not Run} \\
TC-013 & Supported browsers and synthetic Vietnamese names, dates, currencies; keyboard-only user. & Complete core workflows using keyboard; review labels, errors, contrast, date/time, and currency formats. & Record browser/version and findings; no blocking accessibility or localization defect remains at acceptance. & \weak{Not Run} \\
\bottomrule
\end{longtable}
\endgroup

For an executed run, add the actual result, tester, timestamp, build ID, environment, evidence link and, where relevant, defect ID. The test may then be marked Pass, Fail or Blocked. \weak{Until real execution evidence supports those fields, the status stays Not Run.}